% Generated draft: requires mathematical and source review.
\documentclass[11pt]{article}
\usepackage[margin=25mm]{geometry}
\usepackage{fontspec}
\setmainfont{Latin Modern Roman}
\usepackage{amsmath,amssymb,mathtools,unicode-math}
\setmathfont{Latin Modern Math}
\usepackage{enumitem,needspace,xurl,hyperref}
\setlength{\parindent}{0pt}
\setlength{\parskip}{5pt}
\setlength{\emergencystretch}{3em}
\newcommand{\N}{\mathbb{N}}
\newcommand{\Z}{\mathbb{Z}}
\newcommand{\Q}{\mathbb{Q}}
\newcommand{\R}{\mathbb{R}}
\newcommand{\C}{\mathbb{C}}
\newcommand{\F}{\mathbb{F}}
\newcommand{\A}{\mathbb{A}}
\newcommand{\PP}{\mathbb{P}}
\newcommand{\E}{\mathbb{E}}
\newcommand{\eps}{\varepsilon}
\newcommand{\dd}{\,\mathrm{d}}
\begin{document}
\section*{UnsolvedMath ID: 800575}
\textbf{Source problem code:} catalogue entry
\subsection{Definitions and mathematical statement}
Let \(X\) be a topological space. Under the usual nontrivial-cover convention, an open cover of \(X\) is a family \(\mathcal U\) of open subsets of \(X\) such that \(\bigcup\mathcal U=X\) and \(X\notin\mathcal U\). Write \(\mathcal O(X)\) for the class of such covers. An \(\omega\)-cover is an open cover \(\mathcal U\) such that every finite subset of \(X\) is contained in some member of \(\mathcal U\); write \(\Omega(X)\) for this class. A \(\gamma\)-cover is an infinite open cover \(\mathcal U\) such that, for each \(x\in X\), the set \(\{U\in\mathcal U:x\notin U\}\) is finite; write \(\Gamma(X)\) for this class.

For classes \(\mathcal A(X)\) and \(\mathcal B(X)\) of covers, the property \(\mathsf S_{\mathrm{fin}}(\mathcal A,\mathcal B)\) means that for every sequence \((\mathcal U_n)_{n\geq 1}\) with \(\mathcal U_n\in\mathcal A(X)\), there are finite subfamilies \(\mathcal F_n\subseteq\mathcal U_n\) such that
\[
\bigcup_{n\geq 1}\mathcal F_n\in\mathcal B(X).
\]
The property \(\mathsf U_{\mathrm{fin}}(\mathcal A,\mathcal B)\) means that for every sequence \((\mathcal U_n)_{n\geq 1}\) with \(\mathcal U_n\in\mathcal A(X)\) and with no \(\mathcal U_n\) admitting a finite subcover, there are finite subfamilies \(\mathcal F_n\subseteq\mathcal U_n\) such that
\[
\left\{\bigcup\mathcal F_n:n\geq 1\right\}\in\mathcal B(X).
\]
Finite selections are allowed to be empty. The two displayed outputs are families of sets, so repeated members are identified.

\textbf{Unresolved scope and conventions.} The supplied record does not specify the ambient class of spaces, its separation or metrizability assumptions, or whether covers must be countable. The definitions above display the usual nontrivial-cover interpretation; if only countable covers are intended, all three cover classes must be restricted accordingly. These conventions are not confirmed by the supplied excerpt. Let \(\mathscr C\) denote the intended, presently unspecified, class of spaces, with a fixed cover convention.

\textbf{Problem.} Is it true that
\[
\forall X\in\mathscr C,\qquad
X\models\mathsf U_{\mathrm{fin}}(\mathcal O,\Omega)
\quad\Longleftrightarrow\quad
X\models\mathsf S_{\mathrm{fin}}(\Gamma,\Omega)?
\]
Here the equality in the source is interpreted as equality of the classes of spaces satisfying the two properties. If this equivalence fails, is it nevertheless true that
\[
\forall X\in\mathscr C,\qquad
X\models\mathsf U_{\mathrm{fin}}(\mathcal O,\Gamma)
\quad\Longrightarrow\quad
X\models\mathsf S_{\mathrm{fin}}(\Gamma,\Omega)?
\]
The second question retains the source's qualification ``if not.'' A definitive formulation requires recovery of the intended ambient class and cover conventions.
\subsection{Short English statement}
For the intended class of spaces, ask whether \(\mathsf U_{\mathrm{fin}}(\mathcal O,\Omega)\) and \(\mathsf S_{\mathrm{fin}}(\Gamma,\Omega)\) characterize the same spaces. If they do not, ask whether \(\mathsf U_{\mathrm{fin}}(\mathcal O,\Gamma)\) implies \(\mathsf S_{\mathrm{fin}}(\Gamma,\Omega)\). The ambient class and precise cover conventions are absent from the supplied record.
\subsection{Sources}
\begin{itemize}
\item Local UnsolvedMath registry, permanent ID \(800575\), source problem code \texttt{catalogue entry}; title: \emph{Selection Principles in Mathematics --- Problem 1 ([coc2])}. The registry attributes the problem to Tsaban, gives the proposed year as 2006, and labels it open. This is a supplied status label, not an independently verified mathematical or literature assessment.
\item Supplied source list: Tsaban, \emph{Selection Principles and Special Sets of Reals: Open Problems} (2006), Problem 1. Supplied URL: \url{https://arxiv.org/abs/math/0609601}. The record reports access on 2026-10-16 and extraction by \texttt{arxiv\_latex}; these are supplied provenance details, not a claim of access or verification for this draft.
\item The supplied attribution/context is the reference label \texttt{[coc2]}. Its bibliographic entry is not supplied, so its authorship and relationship to the survey attribution remain unresolved.
\end{itemize}
\end{document}
